\medskip
3. \underline{二次超曲面的上同调}

3.1. 设 $p:X\longrightarrow S$ 为 $S$ 上维数为 $n$ 的光滑二次超曲面, $\ell$ 为在 $S$ 上可逆的素数. 以
\begin{equation*}
\eta\in H^0\bigl(S,R^2p_*\mathbb Z_\ell(1)\bigr)
\end{equation*}
表示``超平面截面的上同调类''. 在 $S$ 的平展拓扑下局部地, $P(X)\simeq\mathbb P(V^*)$; 按定义, 此时 $\eta$ 为第一 Chern 类
\begin{equation*}
c_1(\mathcal O(1))\in H^2(X,\mathbb Z_\ell(1)).
\end{equation*}
在 $H^0(S,R^2p_*\mathbb Z_\ell(1))$ 中的像.
\footnote{印本中所列定义写作 $R^1p_*$, 而随后一句和 $c_1(\mathcal O(1))$ 的次数都要求 $R^2p_*$.}

3.2. 假设 $n=2m$. 此时 $X$ 的母空间 $D$ 定义上同调类 $c\ell'(D)\in H^n(X,\mathbb Z_\ell(m))$, 将其在 $H^0(S,R^np_*\mathbb Z_\ell(m))$ 中的像记为 $c\ell(D)$. 由 2.8, 后者只依赖于 $Z(X)$ 的截面 $e(D)$. 因此得到 $S$ 上平展层的态射
\begin{equation}\tag{3.2.1}
c\ell:Z(X)\longrightarrow R^np_*\mathbb Z_\ell(m).
\end{equation}

\underline{定理} 3.3. \underline{设 $p:X\longrightarrow S$ 为 $S$ 上维数为 $n$ 的光滑二次超曲面, $\ell$ 为在 $S$ 上可逆的素数.}

\begin{enumerate}[label=(\roman*),leftmargin=2.6em]
\item $R^{2i+1}p_*\mathbb Z_\ell=0$.

%%%%%%%%%% source scan idx 83 | folio 76 %%%%%%%%%%

\item 若 $0\leq 2i<n$ (分别地, 若 $n<2i\leq 2n$), 则 $R^{2i}p_*\mathbb Z_\ell(i)$ 与常层 $\mathbb Z_\ell$ 典范同构, 并由 $\eta^i$ (分别地, $\eta^i/2$) 生成.

\item 若 $n=2m$, 则从 (3.2.1) 得到的映射
\begin{equation}\tag{3.3.1}
\mathbb Z_\ell^{Z(X)}\longrightarrow R^np_*\mathbb Z_\ell(m),
\end{equation}
为同构. 若 $\alpha$ 和 $\beta$ 为 $Z(X)$ 的两个不相交截面 (假定它同构于 $S\amalg S$, 局部地总是如此), 则
\begin{align*}
\text{(a)}\quad &\eta^m=c\ell(\alpha)+c\ell(\beta),\\[1ex]
\text{(b)}\quad &
\begin{cases}
\operatorname{Tr}(c\ell(\alpha)^2)
=\operatorname{Tr}(c\ell(\beta)^2)=1,
& c\ell(\alpha)c\ell(\beta)=0,\quad m\text{ 为偶数},\\[1ex]
c\ell(\alpha)^2=c\ell(\beta)^2=0,
& \operatorname{Tr}(c\ell(\alpha)c\ell(\beta))=1,
  \quad m\text{ 为奇数},
\end{cases}\\[1ex]
\text{(c)}\quad &\eta\bigl(c\ell(\alpha)-c\ell(\beta)\bigr)=0.
\end{align*}
\end{enumerate}

\footnote{印本在 (iii) 中的映射上重复使用编号 (3.2.1), 但证明本身将此映射引用为 (3.3.1).}

这里只为查考而重述断言 (i) 和 (ii) (XI, 1.6, 在 (i) 中 $2i+1=n$ 时再用 XI, 2.6 补充).

只须在 $S$ 为代数闭域 $k$ 的谱时验证 (iii) (甚至可以取 $S=\Spec(\mathbb C)$).

\underline{公式 (a).} 可假定 $X$ 为 $\mathbb P^{2m+1}(k)$ 中方程为
\begin{equation*}
\sum_{i=0}^{m}x_i x_{i+m+1}=0.
\end{equation*}
的二次超曲面. 类 $\eta^m$ 是由
\begin{equation*}
x_i=0\qquad(0\leq i<m),
\end{equation*}
定义的平面截面的上同调类, 而此圈是由下式定义的母空间 $D_1,D_2$ 之和

%%%%%%%%%% source scan idx 84 | folio 77 %%%%%%%%%%

\begin{align*}
D_1:&\quad x_i=0 &&(0\leq i\leq m),\\
D_2:&\quad x_i=0 &&(0\leq i<m),\qquad x_{2m+1}=0.
\end{align*}
由 1.12, $e(D_1)\neq e(D_2)$, 因而
\begin{equation*}
\eta^m=c\ell(D_1)+c\ell(D_2)
=c\ell(\alpha)+c\ell(\beta).
\end{equation*}

\footnote{印本中的环境射影空间为 $\mathbb P^{m+1}$, 但所列非退化方程有 $2m+2$ 个坐标. 所需的环境空间为 $\mathbb P^{2m+1}$.}

\underline{公式 (b).} 若母空间 $D_1,D_2$ 不相交 (分别地, 在一点处横截相交), 则 $c\ell(D_1)c\ell(D_2)=0$ (分别地, $\operatorname{Tr}(c\ell(D_1)c\ell(D_2))=1$). 由 1.12, 当 $m$ 为偶数时有 $e(D_1)\neq e(D_2)$ (分别地, $e(D_1)=e(D_2)$); 当 $m$ 为奇数时有 $e(D_1)=e(D_2)$ (分别地, $e(D_1)\neq e(D_2)$). 由此得到公式 (b).

\underline{公式 (c).} 当 $n=0$ 时显然成立. 否则, 由 (ii), 只须证明
\begin{equation*}
\eta^m\bigl(c\ell(\alpha)-c\ell(\beta)\bigr)=0.
\end{equation*}
事实上,
\begin{equation*}
\operatorname{Tr}\!\left(
\eta^m(c\ell(\alpha)-c\ell(\beta))\right)
=\operatorname{Tr}\!\left(c\ell(\alpha)^2-c\ell(\beta)^2\right)=0.
\end{equation*}

\underline{验证.} 当 $m$ 为偶数时,
\begin{equation*}
2=\operatorname{Tr}(\eta^{2m})
=\operatorname{Tr}\!\left(c\ell(\alpha)^2
+2c\ell(\alpha)c\ell(\beta)+c\ell(\beta)^2\right)=1+1;
\end{equation*}
当 $m$ 为奇数时,
\begin{equation*}
2=\operatorname{Tr}(\eta^{2m})
=\operatorname{Tr}\!\left(c\ell(\alpha)^2
+2c\ell(\alpha)c\ell(\beta)+c\ell(\beta)^2\right)=2.
\end{equation*}
\footnote{在这两个最高次相交计算中, 印本以 $\eta^2$ 代替 $\eta^{2m}$, 且有一处将 $c\ell(\alpha)^2$ 写作 $c\ell(\alpha^2)$.}

两个二元双线性型
\begin{equation*}
B(x,y)=x_1y_2+x_2y_1,
\qquad
B(x,y)=x_1y_1+x_2y_2
\end{equation*}
的判别式均为 $\pm1$, 由此即得 (3.3.1) 为同构.

%%%%%%%%%% source scan idx 85 | folio 78 %%%%%%%%%%

\underline{验证} 3.4. 设 $X$ 为有限域 $\mathbb F_q$ 上维数为 $n$ 的光滑二次超曲面. 由 3.3, Lefschetz 迹公式给出经典公式
\begin{equation*}
\#X(\mathbb F_q)=
\begin{cases}
\displaystyle\sum_{i=0}^{n}q^i,
  & n\text{ 为奇数},\\[2ex]
\displaystyle\sum_{i=0}^{n}q^i+\varepsilon q^m,
  & n=2m,\quad \varepsilon=\pm1.
\end{cases}
\end{equation*}
当 $n$ 为偶数时, 若 $X$ 分裂, 即具有有理母空间, 则 $\varepsilon=1$, 否则 $\varepsilon=-1$.

3.5. 设 $f:X\longrightarrow S$ 为维数 $n=2m$ 的光滑二次超曲面. $R^nf_*\mathbb Z_\ell(m)$ 的\underline{本原部分}(primitive part)是 $\eta^m$ 的正交补; 沿用 3.3 的记号, 它由 $c\ell(\alpha)-c\ell(\beta)$ 生成. $R^nf_*\mathbb Z_\ell(m)$ 的\underline{本原商}(primitive quotient)是模去 $\mathbb Z_\ell\eta^m$ 所得的商. 将它们与 $\mathbb Z_\ell(-m)$ 作张量积, 即得 $R^nf_*\mathbb Z_\ell$ 的本原部分和本原商.

